\section{Capacité spectrale, mémoire énergétique et transport thermique}
\label{v:sec:capacidad-espectral-termica}

L’application positive entre température et énergie de décision est construite
ici sur le porteur gradué et sa référence marquée. Cette section
conserve l’origine
énergétique qu’une normalisation ne voit pas et compose cette lecture avec le calendrier,
l’aire, la gravitation et les distributions thermiques de l’article. La chaîne
causale demeure
\[
 \mathrm{APP}\longrightarrow\mathrm{TRIT}\longrightarrow\mathrm{TPK}
 \longrightarrow X_{\rm enr}
 \longrightarrow\mathcal C_{\rm cont}^{\rm disc}
 \longrightarrow\text{support thermique}
 \longrightarrow\text{reconnaissance physique}.
\]
Les constantes mathématiques et physiques utilisées sont des sorties HMT--MD ou
des représentations ultérieures de cet état. Aucun chiffre métrologique ne sélectionne
les germes, les secteurs ou le spectre. En particulier, le scalaire
$1{,}380649\times10^{-23}$ s’obtient d’abord comme intensité sans dimension
marquée ; sa composition avec le système de coordonnées énergétiques et thermiques est démontrée
ensuite.

\subsection{Porteur gradué et deux coordonnées conservées}

Soient $\mathscr H_{N,L}$ cinq espaces finis, avec
\begin{equation}
 \dim\mathscr H_{0,0}=1,\quad
 \dim\mathscr H_{1,2}=380,\quad
 \dim\mathscr H_{2,0}=1,\quad
 \dim\mathscr H_{2,1}=24,\quad
 \dim\mathscr H_{2,2}=624.
 \label{v:eq:fibras-termicas-NL}
\end{equation}
Leur somme directe $\mathscr H_{\rm th}$ conserve simultanément le degré
$N$ du préfixe et le degré $L$ de mémoire d’incidence. Le polynôme de
Hilbert bivarié est
\begin{equation}
 Z(u,v)=1+380uv^2+u^2(1+24v+624v^2).
 \label{v:eq:polinomio-graduado-termico}
\end{equation}
Par conséquent,
\[
 Z(u,1)=1+380u+649u^2,\qquad
 Z(1,v)=2+24v+1004v^2.
\]
L’égalité $649=1+24+624$ décrit les fibres de $L$ au degré $N=2$ ;
l’identité indépendante $649=1+2\cdot324$ décrit la décomposition
par involution du porteur. Ces deux significations ne sont pas échangées.

Les dimensions $380$ et $649$ proviennent de deux opérations distinctes. La
première est la partie non diagonale de l’endomorphisme sur le porteur de
capacité $20$ ; la seconde réunit la direction scalaire et deux orientations
du secteur de $324$ chemins. Pour expliciter cette seconde construction,
soit
\[
 X_\times=(\mathbb Z/9\mathbb Z)^2\times\{N,E,S,O\},
 \qquad |X_\times|=324,
\]
et soit $\mathfrak e:X_\times\to\mathcal S_\times$ l’émetteur à six
fenêtres, avec $26$ signatures. Ses fibres ont pour cardinal $8$ pour $24$ signatures,
$24$ pour $555555$ et $108$ pour $000000$. Avec
$I_{24}=\{1,\ldots,24\}$, on construit l’arbre microscopique
\begin{equation}
 \widetilde{\mathcal T}=\{\Omega\}\sqcup\{a_i:i\in I_{24}\}
 \sqcup\{b_{i,y}:i\in I_{24},\ y\in X_\times\},
 \qquad |\widetilde{\mathcal T}|=7801.
 \label{v:eq:arbol-microscopico-7801}
\end{equation}
La lecture $p(b_{i,y})=b_{i,\mathfrak e(y)}$, qui fixe $\Omega$ et les
$a_i$, a pour image $\mathcal T$ de cardinal
$1+24+24\cdot26=649$. Si
$m_s=|\mathfrak e^{-1}(s)|$, la mesure
\begin{equation}
 \nu_{\rm mic}(\Omega)=\nu_{\rm mic}(a_i)=1,
 \qquad \nu_{\rm mic}(b_{i,y})=m_{\mathfrak e(y)}^{-1}
 \label{v:eq:medida-arbol-microscopico}
\end{equation}
donne une masse un à chaque fibre de $p$. Les taux sont $r>0$ dans les deux sens
entre $\Omega$ et $a_i$, $r/m_{\mathfrak e(y)}$ de $a_i$ vers $b_{i,y}$
et $r$ dans le sens inverse.

Sur $\ell^2(\mathcal T)$ avec mesure de comptage et
$\ell^2(\widetilde{\mathcal T},\nu_{\rm mic})$, le relèvement
\begin{equation}
 J_p g=g\circ p
 \label{v:eq:levantamiento-fibras-Jp}
\end{equation}
est une isométrie de la classe ou signature vers l’arbre microscopique pondéré.
En effet,
\[
 \|J_pg\|_{\nu_{\rm mic}}^2
 =\sum_{z\in\mathcal T}|g(z)|^2
   \sum_{x:p(x)=z}\nu_{\rm mic}(x)=\|g\|^2,
 \qquad J_p^*J_p=I.
\]
L’adjoint est la moyenne pondérée
\begin{equation}
 (J_p^*f)(z)=\sum_{x:p(x)=z}\nu_{\rm mic}(x)f(x),
 \label{v:eq:adjunto-fibras-Jp}
\end{equation}
non le relèvement $J_p$ : ils ont des domaines opposés et l’un copie une fonction
de classe tandis que l’autre calcule la moyenne sur sa fibre. Le bilan détaillé
$\nu_{\rm mic}(a_i)r/m_{\mathfrak e(y)}
=\nu_{\rm mic}(b_{i,y})r$ et les sommes des taux prouvent, pour les
générateurs microscopique et agrégé,
$\widetilde L J_p=J_pL_{\mathcal T}$. Le bilan à l’intérieur d’une fibre ne
détermine pas à lui seul les poids entre degrés ; ceux-ci sont fixés par les
courants de bord construits plus bas.

La variable
\begin{equation}
 Q:=N-L
 \label{v:eq:carga-memoria-Q}
\end{equation}
sépare deux états qu’une marginale peut identifier. Par exemple, connaître
$Z(u,1)$ ne récupère pas la distribution de $L$. Les cinq cellules de
\eqref{v:eq:fibras-termicas-NL} et leurs multiplicités reconstruisent le
porteur conjoint ; elles ne reconstruisent pas les cohérences qu’une compression a
éliminées. C’est le domaine exact de la lecture.

\subsection{Relèvement positif de la mémoire énergétique}

Dans les cinq cellules, $Q$ prend les valeurs $0,-1,2,1,0$. Introduisons un
registre de mémoire $R$ avec niveaux $0,1,2,3$ et l’isométrie
\begin{equation}
 V_{\rm mem}:\mathscr H_{N,L}\longrightarrow
 \mathscr H_{N,L}\otimes\mathbb C^4,\qquad
 V_{\rm mem}\xi=\xi\otimes e_{N-L+1}.
 \label{v:eq:isometria-memoria-positiva}
\end{equation}
Le décalage d’une unité transforme la charge signée en un niveau non
négatif sans la perdre. Si $R e_j=j e_j$, alors
\begin{equation}
 V_{\rm mem}^*(L\otimes I+I\otimes R)V_{\rm mem}=N+I.
 \label{v:eq:balance-positivo-memoria}
\end{equation}
Dans chaque cellule, cette identité est simplement
$L+(N-L+1)=N+1$, mais son écriture opératorielle conserve les dégénérescences,
les superpositions à l’intérieur de chaque fibre et toute purification sur un
système auxiliaire.

Pour une échelle $\epsilon_0>0$, l’égalité
\begin{equation}
 V_{\rm mem}^*\epsilon_0(L\otimes I+I\otimes R)V_{\rm mem}
 =\epsilon_0(N+I)
 \label{v:eq:energia-lift-memoria}
\end{equation}
démontre la conservation exacte de l’énergie dans le relèvement. Le filtre qui
retient une branche doit aussi conserver le marqueur de succès ; sinon,
la densité conditionnelle masque le coût de préparation. Le relèvement est
réversible sur son image et a un coût logique nul ; éliminer le registre
$R$ est une opération distincte, soumise au bilan de Landauer de la section
précédente.

\subsection{Spectre décimal produit par le lacet de mémoire}

La décomposition d’incidence $8{:}1$ détermine le couplage
\[
 \mathcal Q=\frac13
 \begin{pmatrix}\sqrt8 I&-I\\ I&\sqrt8 I\end{pmatrix}.
\]
Le lacet orienté de retour est
\[
 H_1=\begin{pmatrix}2&1\\1&1\end{pmatrix},\qquad
 H_1H_1^T=\begin{pmatrix}5&3\\3&2\end{pmatrix}.
\]
Dans la préparation gaussienne positive déclarée, la covariance visible
normalisée est
\begin{equation}
 W=\frac{8I+H_1H_1^T}{9},\qquad
 \det W=\frac{121}{81},\qquad
 \nu_W=\sqrt{\det W}=\frac{11}{9}.
 \label{v:eq:williamson-decimal}
\end{equation}
Le théorème spectral à un mode donne alors
\begin{equation}
 r=\frac{\nu_W-1}{\nu_W+1}=\frac1{10},\qquad
 \rho_r=(1-r)\sum_{j\ge0}r^j\,|j\rangle\langle j|.
 \label{v:eq:estado-geometrico-decimal}
\end{equation}
La raison décimale s’obtient à partir de la covariance du lacet ; elle n’est pas introduite comme
base numérique cible. Le symbole $\nu_W$ est la valeur propre symplectique de
Williamson et est distinct de la mesure microscopique $\nu_{\rm mic}$ de
\eqref{v:eq:medida-arbol-microscopico}. Le contrôle $H_0=I$ produit $r=0$,
de sorte que le spectre dépend effectivement de l’holonomie choisie.

La division euclidienne $j=3n+b$, $0\le b<3$, définit l’unitaire
\begin{equation}
 U_3|j\rangle=|n\rangle\otimes|b\rangle,\qquad
 U_3\rho_rU_3^*=\rho_{r^3}\otimes
 \frac{\operatorname{diag}(1,r,r^2)}{1+r+r^2}.
 \label{v:eq:agrupamiento-tritico-espectral}
\end{equation}
La preuve est terme à terme :
$(1-r)r^{3n+b}=(1-r^3)(r^3)^n r^b/(1+r+r^2)$.
Ainsi,
\begin{equation}
 q=r^3=10^{-3},\qquad
 \sigma_r=\frac{\operatorname{diag}(100,10,1)}{111}.
 \label{v:eq:razon-milenaria-y-residuo}
\end{equation}
Le résidu conserve les trois positions intermédiaires. Si $S|j\rangle=|j+1\rangle$,
son transport est
\[
 U_3SU_3^*=I\otimes(|1\rangle\langle0|+|2\rangle\langle1|)
 +S_n\otimes|0\rangle\langle2|,
\]
de sorte que la troisième avancée fait revenir le résidu et augmente le quotient. Le retour
de phase et le retour de l’état ne sont pas confondus.

Avec $H_W=\epsilon_a(N+\tfrac12)$,
\begin{equation}
 U_3H_WU_3^*
 =\epsilon_a\left(3N_n\otimes I+I\otimes R+\tfrac12I\right).
 \label{v:eq:hamiltoniano-agrupado-decimal}
\end{equation}
La séparation du facteur de quotient est $3\epsilon_a$, tandis que la
température reste fixe :
\[
 e^{-\beta_W\epsilon_a}=10^{-1},\qquad
 e^{-\beta_W3\epsilon_a}=10^{-3},\qquad
 k_BT_W=\frac{\epsilon_a}{\log10}.
\]

\subsection{Section marquée et évaluation complète}

Le renouvellement de~\eqref{v:eq:tupla-renovacion-capacidad} fournit
l’origine $23$ et le raffinement par trois fournit $23,26,29$. Avec
$\lambda_j/\lambda_0=r^j$ et les multiplicités du porteur,
\begin{equation}
 \boxed{
 \mathcal B
 =\frac{\lambda_{23}+380\lambda_{26}+649\lambda_{29}}{\lambda_0}
 =10^{-23}\left(1+\frac{380}{1000}+
                    \frac{649}{1000^2}\right)
 =1{,}380649\times10^{-23}.}
 \label{v:eq:evaluacion-termica-completa}
\end{equation}
Le résidu $b=2$ isolé sélectionnerait $2,5,8,\ldots$ et ne détermine pas le
seuil $23$ ; celui-ci provient du temps de renouvellement et de l’indice de capacité.
La mantisse est conservée dans la densité normalisée au moyen du marqueur du
secteur ; le poids de sélection conserve en outre $10^{-23}$. Ainsi, ni la
mantisse ni l’exposant ne sont ajoutés à partir d’une table externe.

Soit
$\mathcal D=\bigoplus_{n=0}^2\mathcal D_n$, avec
$(d_0,d_1,d_2)=(1,380,649)$,
$\dim\mathcal D_n=d_n$ et $N|_{\mathcal D_n}=nI$. Nous définissons l’opérateur
de capacité et la référence spectrale
\begin{equation}
 C=23I+3N,\qquad
 \eta=\bigoplus_{n=0}^2r^{23+3n}I_{\mathcal D_n},\qquad
 b_*:=\operatorname{Tr}\eta=\mathcal B.
 \label{v:eq:referencia-termica-eta}
\end{equation}
La référence $\eta$ conserve séparément les degrés $23,26,29$ et
reste fixe lorsqu’un état d’essai varie. Pour
$x=e^{-\epsilon_0/\tau}$, la partition relative et son état sont
\begin{equation}
 Z_N(x)=1+380x+649x^2,\qquad
 \rho_N(x)=\frac{x^N}{Z_N(x)},\qquad
 p_0=\frac1{Z_N(x)}.
 \label{v:eq:particion-portador-termico}
\end{equation}
Le marqueur $p_0$ récupère $Z_N$ même après normalisation. Dans
$q=10^{-3}$, $Z_N(q)=1{,}380649$,
$\eta=10^{-23}q^N$ et $\eta/b_*=\rho_N(q)$ ; l’origine marquée produit
\eqref{v:eq:evaluacion-termica-completa}. L’énergie libre et le degré moyen sont
\begin{equation}
 F_N(\tau)=-\tau\log Z_N(x),\qquad
 \langle N\rangle_x=x\frac{d}{dx}\log Z_N(x)
 =\frac{380x+1298x^2}{Z_N(x)}.
 \label{v:eq:energia-libre-grado-medio}
\end{equation}

\subsection{Courants de bord et rigidité de l’origine sectorielle}

L’incidence surfacique fournit $18+18$ liens tritiques. Soient
$u=N_{90}$, $v=N_{120}$, $a=u+v$ et
\[
 F=3u+4v,\qquad H_R=\epsilon_0F,
 \quad \epsilon_0=30\Delta\tau_R>0,
 \quad \tau_R=\frac{\hbar c}{2\pi R}.
\]
La provenance demeure dans $(u,v)$ : la même aire peut correspondre à
des énergies différentes. Dans un lien,
$T^+=|1\rangle\langle0|+|2\rangle\langle1|$ et $T^-=(T^+)^*$ satisfont
$[n,T^\pm]=\pm T^\pm$. Pour $e,f=1,\ldots,18$, nous définissons
\[
 J_{ef}=T^+_{120,f}T^-_{90,e},\qquad
 B_c=\left(\sum_eT^+_{90,e}\right)^3
     \left(\sum_fT^-_{120,f}\right)^2.
\]
La règle $[F,AB]=[F,A]B+A[F,B]$ prouve
\[
 [F,J_{ef}]=J_{ef},\quad[a,J_{ef}]=0,
 \qquad [F,B_c]=B_c,\quad[a,B_c]=B_c.
\]
La suite des occupations
$(0,2)\to(3,0)\to(2,1)$ a pour énergies primitives $8,9,10$ et pour aires
$2,3,3$ : le même incrément énergétique conserve deux actions surfaciques
distinguables.

Les coefficients $c_j=[z^j](1+z+z^2)^{18}$ comptent les occupations de
chaque canal : $c_0=1,c_1=18,c_2=171,c_3=1122$. Le caractère de frontière
\[
 P(x)=(1+x^3+x^6)^{18}(1+x^4+x^8)^{18}
\]
donne $g_7=324,g_8=171,g_9=1122,g_{10}=3078$. Pour les rangs
$(1,380,649)$, le premier triplet admissible d’énergies consécutives commence
en $8$. Les débuts inférieurs sont exclus parce que, en utilisant
$g_0=1$, $g_1=g_2=g_5=0$, $g_3=g_4=18$, $g_6=171$, $g_7=324$ et
$g_8=171$, chaque triplet précédent ne satisfait pas au moins l’un des rangs.

La compression connexe est construite en choisissant un état $(0,2)$,
$380$ états $(3,0)$ et $649$ voisins distincts de ceux-ci au moyen des
$J_{ef}$. Chaque occupation de somme trois possède au moins deux prédécesseurs de
somme deux et chaque prédécesseur possède au plus $18$ successeurs. Par conséquent,
les $380$ états produisent au moins
$\lceil760/18\rceil=43$ prédécesseurs distincts et
$18\cdot43=774$ voisins $(2,1)$, parmi lesquels on peut en choisir $649$.
Le courant $B_c$ relie l’état initial aux $380$ états
intermédiaires et chaque voisin final conserve un antécédent sélectionné ; le
graphe résultant est connexe. Ainsi sont construits, et non simplement
postulés, le sous-espace comprimé, ses étiquettes et ses courants.

La compression de bord réalise les degrés $N=0,1,2$ aux énergies
$8\epsilon_0,9\epsilon_0,10\epsilon_0$. Soient
$\mathcal D_n=\ker(N-nI)$,
$P_F=\mathbf1_{\{N=1\}}$,
$B_N:\mathcal D_0\to\mathcal D_1$ et
$J_N:\mathcal D_1\to\mathcal D_2$ les courants comprimés non nuls ;
leurs adjoints parcourent les arêtes dans le sens descendant. Alors
\begin{equation}
 [N,B_N]=B_N,\quad [N,J_N]=J_N,\qquad
 [P_F,B_N]=B_N,\quad [P_F,J_N]=-J_N.
 \label{v:eq:conmutadores-corrientes-termicas}
\end{equation}
Pour la correction diagonale affine
\[
 H_\delta=\epsilon_0(8I+N)+\delta P_F
\]
il découle des commutateurs que
\[
 [H_\delta,B_N]=(\epsilon_0+\delta)B_N,
 \qquad
 [H_\delta,J_N]=(\epsilon_0-\delta)J_N.
\]
Exiger que les deux courants conservent le saut de bord $\epsilon_0$, ou
simplement que les deux sauts soient égaux, impose
\begin{equation}
 \delta=0.
 \label{v:eq:rigidez-desfase-corrientes}
\end{equation}
La conclusion classe ce déphasage sectoriel ; elle n’affirme pas que l’on a
classé tout le commutant du Hamiltonien.

En attribuant à chaque arête ascendante le taux $xr_e$ et à l’arête descendante le taux
$r_e>0$, avec le même $r_e$ dans les deux sens, on obtient le bilan détaillé
par rapport aux poids $x^N$. La connexité rend unique la distribution
stationnaire sur le porteur comprimé et récupère
\eqref{v:eq:particion-portador-termico}. Cette dynamique est un semi-groupe de
Markov de préparation thermique, non une évolution unitaire isolée.

\subsection{Équivalence thermique avec origine et marqueur}

Soient $(\mathscr H_i,H_i,\tau_i,e_i,P_i)$ deux réalisations, avec
$\tau_i>0$, l’origine $e_i$, le projecteur de support $P_i$ et un isomorphisme
isométrique $J:P_1\mathscr H_1\to P_2\mathscr H_2$. L’équivalence thermique
pertinente est
\begin{equation}
 J\frac{H_1-e_1I}{\tau_1}
 =\frac{H_2-e_2I}{\tau_2}J,\qquad
 JP_1=P_2J.
 \label{v:eq:equivalencia-termica-centrada}
\end{equation}
Elle n’exige ni $H_1=H_2$ ni $\tau_1=\tau_2$. Elle implique
$Jf((H_1-e_1I)/\tau_1)=f((H_2-e_2I)/\tau_2)J$ pour toute fonction spectrale
bornée $f$, et transporte les états de Gibbs centrés. Pour les observables
qui dépendent de l’énergie absolue, il faut transporter en outre $e_i$ et
l’échelle : normaliser n’autorise pas à les effacer.

Dans la compression précédente, il existe une isométrie basée $\iota_N$ telle que
\[
 H_R\iota_N=\iota_N\epsilon_0(8I+N).
\]
Le courant $J_N:\mathcal D_1\to\mathcal D_2$ et l’isométrie
$\iota_N$ sont des opérateurs distincts. Pour la lecture chronologique, on fixe
\begin{equation}
 C=23I+3N,\qquad
 H_{\rm cap}=\frac{\epsilon_a\log10}{\log3}C,
 \qquad \tau_{\rm clk}=\frac{\epsilon_a}{\log3}.
 \label{v:eq:lector-capacidad-escalas}
\end{equation}
La section relative de Williamson utilise, en revanche,
$H_{W,{\rm rel}}=\epsilon_a C$ et
$\tau_W=\epsilon_a/\log10$. Par conséquent,
$H_{\rm cap}/\tau_{\rm clk}=H_{W,{\rm rel}}/\tau_W=C\log10$,
bien que leurs échelles d’énergie et de température soient différentes.

Avec $\tau_*=\epsilon_0/\log1000$, la comparaison centrée est
\begin{equation}
 \frac{H_R-8\epsilon_0I}{\tau_*}\iota_N
 =\iota_N\frac{H_{\rm cap}-23\tau_{\rm clk}\log10\,I}
                    {\tau_{\rm clk}}
 =\iota_N(\log1000)N.
 \label{v:eq:transporte-termico-centrado}
\end{equation}
L’égalité conserve les degrés, les énergies, le marqueur et l’origine sans identifier
$\epsilon_0$ à $\epsilon_a$ : la première échelle appartient à la compression
de bord et la seconde à la section chronologique. Le facteur $1/10$ entre
les traces non centrées et le décalage $\log10$ sont des données de la section,
non des défauts de la densité.

\subsection{Aire, provenance et coût thermique}

Dans les deux canaux surfaciques, soient $N_{90},N_{120}\ge0$ et
\begin{equation}
 a=N_{90}+N_{120},\qquad
 D=90N_{90}+120N_{120}.
 \label{v:eq:ocupacion-procedencia-90-120}
\end{equation}
On récupère exactement
\begin{equation}
 N_{90}=4a-\frac D{30},\qquad
 N_{120}=\frac D{30}-3a.
 \label{v:eq:inversion-procedencia-90-120}
\end{equation}
L’aire peut lire $a$, tandis que le Hamiltonien sectoriel conserve la
provenance :
\begin{equation}
 H_R=\tau_R\Delta D
 =\tau_R\Delta(90N_{90}+120N_{120}).
 \label{v:eq:hamiltoniano-area-procedencia}
\end{equation}
Deux états ayant le même $a$ et une répartition différente entre canaux ont la
même lecture d’occupation totale et un coût distinct. La réduction à un canal,
sa purification et l’identité d’énergie libre finie s’appliquent à l’état
réduit déjà construit ; l’entropie relative quantifie le déficit par rapport
à l’état de Gibbs de référence. Cette composition explique pourquoi
l’origine et la provenance restent nécessaires lorsqu’on relie l’aire,
l’intrication et la fonctionnelle de Barbero--Immirzi.

\subsection{Composition avec l’action, la gravitation et l’électron}

La longueur $L_\alpha$ et l’action orientée $\hbar_a$ fixent
\begin{equation}
 \epsilon_a=\frac{\hbar_ac}{L_\alpha},\qquad
 G_a=\frac{c^3L_\alpha^2}{\hbar_a},\qquad
 k_B\Theta_{{\rm clk},a}\log3=\epsilon_a.
 \label{v:eq:triple-accion-gravedad-temperatura}
\end{equation}
Multiplier les trois identités élimine la section d’action et donne
\begin{equation}
 \boxed{G_a k_B\Theta_{{\rm clk},a}\log3=c^4L_\alpha.}
 \label{v:eq:identidad-gravedad-temperatura}
\end{equation}
L’identité s’obtient dans la même orientation $a$ ; $G_a$ n’est pas introduit
à partir de la métrologie. La lettre $L_\alpha$ désigne la longueur et non le degré $L$ de
\eqref{v:eq:fibras-termicas-NL}.

Si $R_{\rm act}=\hbar_{\rm pre}/\hbar_{\rm ret}$ est le rapport d’action du
retour électronique, la même composition produit
\begin{equation}
 R_{\rm act}
 =\frac{\Theta_{{\rm clk,pre}}}{\Theta_{{\rm clk,ret}}}
 =\frac{G_{\rm ret}}{G_{\rm pre}}.
 \label{v:eq:retorno-accion-termico-gravitatorio}
\end{equation}
Cette égalité transporte la section thermique ; elle n’utilise pas $k_B$ pour sélectionner
le chemin et ne redérive pas dans cet article la loi des masses. Pour une énergie
électronique $E_e$ déjà individualisée, $y_e=E_e/\epsilon_{\rm ret}$ satisfait
\[
 \beta_{\rm clk}E_e=(\log3)y_e,\qquad
 \frac{G_{\rm ret}m_e^2}{\hbar_{\rm ret}c}=y_e^2.
\]
L’origine de l’énergie affecte la masse et le rayon même si elle s’annule dans les probabilités
normalisées.

\subsection{Pont thermométrique marqué et identité effective}

La référence $\eta$ de~\eqref{v:eq:referencia-termica-eta} doit rester
séparée de l’état sur lequel l’énergie est mesurée. Soit $\mathcal H$ un espace
fini, $\rho$ une matrice densité sur $\mathcal H$ et
$H=H^*$ un Hamiltonien à origine enregistrée. Dans
$\widehat{\mathcal D}=\mathcal D\oplus\mathbb C f$, nous définissons
\begin{equation}
 \widehat\eta=\eta\oplus(1-b_*)|f\rangle\langle f|,
 \qquad \Omega_\rho=\widehat\eta\otimes\rho,
 \label{v:eq:referencia-marcada-normalizada}
\end{equation}
et nous notons $P$ le projecteur sur le terme de somme directe $\mathcal D$, étendu
par l’identité de $\mathcal H$. Avec
$s(\rho)=-\operatorname{Tr}(\rho\log\rho)$ et les logarithmes restreints
au support, les deux lectures locales sont
\begin{align}
 \mathcal S_P(\rho)
 &=-\operatorname{Tr}\!\left[P\Omega_\rho
                 (I\otimes\log\rho)\right]
   =b_*s(\rho),\label{v:eq:lector-entropico-marcado}\\
 \mathcal E_P(\rho)
 &=\frac{\operatorname{Tr}\!\left[P\Omega_\rho(I\otimes H)\right]}
         {\operatorname{Tr}(P\Omega_\rho)}
   =\operatorname{Tr}(\rho H).
 \label{v:eq:lector-energetico-condicionado}
\end{align}
En effet, $P\Omega_\rho=\eta\otimes\rho$ et
$\operatorname{Tr}(P\Omega_\rho)=b_*$. La factorisation de la trace donne
$b_*s(\rho)$ et $b_*\operatorname{Tr}(\rho H)$ dans les deux numérateurs ;
diviser le second par $b_*>0$ prouve la lecture énergétique. Ainsi sont
conservées simultanément la référence fixe, l’information relative à
sa marque et l’énergie par état conditionné.

Nous travaillons dans les bases positives transportées des droites d’énergie,
de température et d’entropie, avec $u_{\rm ent}=u_E/u_\Theta$. Dans les formules de
coordonnées, $H$ exprime l’énergie en $u_E$ et $\mathcal S_P$ exprime l’entropie
en $u_{\rm ent}$. La règle de réalisation se compose de quatre opérations :
recevoir la référence fixe~\eqref{v:eq:referencia-marcada-normalizada} ;
évaluer~\eqref{v:eq:lector-entropico-marcado} ;
évaluer~\eqref{v:eq:lector-energetico-condicionado} ; et définir la section
thermique comme conjuguée de ces deux fonctionnelles. Les bases sont transportées avec
les opérateurs, sans introduire de température cible.

\begin{theorem}[Dérivée thermométrique de la lecture marquée]
\label{v:thm:derivada-termometrica-marcada}
Si $H$ n’est pas scalaire et
$\rho_\beta=e^{-\beta H}/Z(\beta)$ pour $\beta>0$, alors
\begin{equation}
 \frac{d\mathcal S_P}{d\mathcal E_P}=b_*\beta,
 \qquad \Theta_P=\frac1{b_*\beta},
 \qquad B_*(\Theta_Pu_\Theta)=\beta^{-1}u_E,
 \label{v:eq:derivada-termometrica-marcada}
\end{equation}
où $B_*(u_\Theta)=b_*u_E$ est l’unique application linéaire positive ayant cette
image de la base.
\end{theorem}
\begin{proof}
Pour $E(\beta)=\operatorname{Tr}(\rho_\beta H)$, la somme spectrale finie
donne $E'=-\operatorname{Var}_{\rho_\beta}(H)<0$ et
$(\log Z)'=-E$. De
$s(\rho_\beta)=\beta E+\log Z$ résulte $s'=\beta E'$. Comme $\eta$
reste fixe,
$\mathcal S_P'=b_*s'$ et $\mathcal E_P'=E'$. Leur quotient et son inverse
prouvent les deux premières identités ; l’image d’une base positive fixe
une unique application linéaire et prouve la troisième.
\end{proof}

Le facteur dépend de la paire de lectures. Si l’énergie et l’entropie sont prises
toutes deux marquées, $(b_*E,b_*s)$, ou toutes deux conditionnées, $(E,s)$, la dérivée
est $\beta$ ; pour la paire marquée--conditionnée
\eqref{v:eq:lector-entropico-marcado}--
\eqref{v:eq:lector-energetico-condicionado}, elle est $b_*\beta$.

La même règle possède une caractérisation variationnelle. Pour $\Theta>0$,
\begin{equation}
 \Phi_\Theta(\rho)=\mathcal E_P(\rho)-\Theta\mathcal S_P(\rho)
 =\operatorname{Tr}(\rho H)-b_*\Theta s(\rho)
 \label{v:eq:funcional-variacional-marcado}
\end{equation}
a l’unique minimiseur
\begin{equation}
 \rho_\Theta^*=\frac{e^{-H/(b_*\Theta)}}
                    {\operatorname{Tr}e^{-H/(b_*\Theta)}}.
 \label{v:eq:minimo-variacional-marcado}
\end{equation}
En effet,
\[
 \Phi_\Theta(\rho)-\Phi_\Theta(\rho_\Theta^*)
 =b_*\Theta D_{\rm rel}(\rho\Vert\rho_\Theta^*)\ge0,
\]
avec
$D_{\rm rel}(\rho\Vert\sigma)
=\operatorname{Tr}[\rho(\log\rho-\log\sigma)]$ ; l’égalité dans la
positivité de l’entropie relative a lieu exactement
pour $\rho=\rho_\Theta^*$. Cela démontre l’existence et l’unicité sans
introduire de température cible.

Nous appliquons maintenant le résultat à l’horloge de capacité. L’orientation
et $\epsilon_a$ étant fixés, soient $H_{{\rm clk},a}=H_{\rm cap}$ et
$\beta_{{\rm clk},a}=\log3/\epsilon_a$. Ce Hamiltonien est positif,
autoadjoint et non scalaire, car ses trois valeurs propres distinctes sont
proportionnelles à $23,26,29$. Alors
\begin{align}
 \beta_{{\rm clk},a}H_{{\rm clk},a}
   &=C\log10,\nonumber\\
 e^{-\beta_{{\rm clk},a}H_{{\rm clk},a}}
   &=10^{-23}q^N=\eta,\qquad q=10^{-3},\nonumber\\
 \rho_{{\rm clk},a}
   &=\frac{q^N}{Z_N(q)}=\frac{\eta}{b_*}.
 \label{v:eq:estado-reloj-capacidad}
\end{align}
La référence $\eta$, indépendante de $\epsilon_a$, reste fixe dans le
premier facteur, tandis que $\rho_\beta$ varie dans une seconde copie de
$\mathcal D$ avec $H_{{\rm clk},a}$ fixe. Évaluer ensuite en
$\beta=\beta_{{\rm clk},a}$ donne
\begin{equation}
 \Theta_{{\rm clk},P,a}=\frac{\epsilon_a}{b_*\log3},
 \qquad
 \boxed{B_*(\Theta_{{\rm clk},P,a}u_\Theta)
 =\frac{\epsilon_a}{\log3}u_E
 =\frac{h_a}{108t_0\log3}u_E.}
 \label{v:eq:identidad-transductor-numeral}
\end{equation}
La dernière égalité utilise la coordonnée d’action déjà fixée
$\epsilon_a=h_a/(108t_0)$ ; elle n’utilise pas la valeur de $b_*$ comme ajustement. Elle n’identifie pas non plus
$\epsilon_a$ à l’échelle de bord $\epsilon_0$.

Si $B_{\rm clk}$ est le transducteur chronologique et que $J_E,J_\Theta$ transportent
respectivement l’énergie et les deux sections thermométriques du même état
de référence, l’unicité entre droites donne le carré typé
\begin{equation}
 B_*J_\Theta=J_EB_{\rm clk}.
 \label{v:eq:cuadrado-transductor-marcado}
\end{equation}
Les deux compositions coïncident sur la section positive génératrice par
\eqref{v:eq:identidad-transductor-numeral} ; la linéarité étend
l’égalité à toute la droite. C’est le pont entre l’évaluation spectrale
de~\eqref{v:eq:evaluacion-termica-completa} et le transducteur dimensionnel :
les domaines, la référence, les lecteurs et la section positive déterminent l’application avant
son évaluation numérique.

\subsection{Conséquences en occupation, rayonnement et information}

Le porteur fini apporte à la théorie de Gibbs la partition
$Z_N(x)$, l’état $\rho_N(x)$, le marqueur $p_0$ et l’énergie libre
\eqref{v:eq:energia-libre-grado-medio}. Les complétions symétrique et
extérieure développées plus loin agissent sur des modes complets et produisent
les lois de Bose--Einstein et de Fermi--Dirac ; le porteur $1/380/649$ ne les
remplace pas. Il leur fournit une section énergétique et une préparation à origine
enregistrée. Dans le régime dilué, la même coordonnée
$x=e^{-\beta(E-\mu)}$ entre dans la limite de Maxwell--Boltzmann déjà démontrée.

Dans la réalisation radiative, la composition
$\beta=(k_B\Theta)^{-1}$ et
$k_B\Theta_{{\rm clk},a}\log3=\epsilon_a$ substitue l’échelle du calendrier
dans la densité de Planck sans modifier le dénombrement des deux polarisations ni sa
dispersion. Les moments conservent les déductions de Stefan--Boltzmann et de
Wien ; la nouvelle section détermine d’où provient l’échelle thermique et quel
marqueur doit être conservé lors du conditionnement d’une branche. Dans l’aire et
l’intrication,~\eqref{v:eq:inversion-procedencia-90-120} empêche que
l’occupation totale efface le coût de provenance. Ainsi, l’occupation, le rayonnement,
l’énergie libre, l’aire et la gravitation reçoivent la même information thermique sans
confondre leurs observables.
